\documentclass[12pt]{article}
\usepackage[T1]{fontenc}
\usepackage[utf8]{inputenc}
\usepackage{lmodern}
\usepackage{textcomp}
\usepackage{enumitem}
\usepackage{geometry}
\geometry{margin=1in}
\begin{document}
\begin{center}
Answer Key - Religious Studies - K-12 Level Assessment 5 \
\small (Answers are ordered exactly as in the question paper.)
\end{center}

\section*{Multiple Choice}
\begin{enumerate}[label=\arabic*.]
\item \textbf{Answer:} C
\\ \textit{Options:} A) The Battle of Uhud; B) The Battle of Camel; C) The Battle of the Trench; D) The Battle of Badr
\\ \textit{Note:} The Battle of the Trench, also known as the Battle of the Confederates, was a pivotal moment in early Islamic history that showcased the strength and resilience of the Muslim community in Medina, ultimately leading to the weakening of their opposition and paving the way for the eventual peaceful conversion of Mecca to Islam.
\item \textbf{Answer:} C
\\ \textit{Options:} A) Priests; B) Servants; C) Rulers, warriors; D) Merchants
\\ \textit{Note:} The kshatriyas are recognized in the traditional Indian caste system as the class responsible for governance and military defense, which is why they are identified as rulers and warriors.
\item \textbf{Answer:} A
\\ \textit{Options:} A) Muhammad's miraculous ascent to heaven; B) Muhammad's migration to Mecca; C) Muhammad's first community in Mecca; D) Muhammad's revelations of the Qur'an
\\ \textit{Note:} The mi'raj refers to the event in Islamic tradition where the Prophet Muhammad is believed to have been miraculously taken from Earth to the heavens, making "Muhammad's miraculous ascent to heaven" the accurate definition of this significant religious experience.
\item \textbf{Answer:} C
\\ \textit{Options:} A) Peter, John, Mark; B) Matthew, Mark, John; C) Matthew, Mark, Luke; D) John, Luke, Mark
\\ \textit{Note:} The synoptic Gospels-Matthew, Mark, and Luke-are called "synoptic" because they share similar stories, teachings, and viewpoints about the life and ministry of Jesus, allowing them to be compared and contrasted in a unified manner.
\item \textbf{Answer:} B
\\ \textit{Options:} A) Kami; B) Kokugaku; C) Samurai; D) Gomadaki
\\ \textit{Note:} Kokugaku, or "National Learning," was a movement in Edo-period Japan that sought to revive and promote indigenous Japanese cultural practices and beliefs while actively rejecting the influences of Confucianism and Buddhism.
\end{enumerate}

\section*{True / False}
\begin{enumerate}[label=\arabic*.]
\setcounter{enumi}{5}
\item \textbf{Answer:} True
\\ \textit{Options:} A) True; B) False
\\ \textit{Note:} Sarasvati is indeed recognized in Hinduism as the goddess of speech and knowledge, and she is traditionally regarded as the mother of the Vedas, which are the sacred texts of Hindu philosophy and spirituality.
\item \textbf{Answer:} True
\\ \textit{Options:} A) True; B) False
\\ \textit{Note:} In Jainism, spiritual liberation, or Moksha, is achieved by freeing oneself from the cycle of Samsara, which is the continuous cycle of birth, death, and rebirth, through right knowledge, right faith, and right conduct.
\item \textbf{Answer:} True
\\ \textit{Options:} A) True; B) False
\\ \textit{Note:} The statement is true because Antiochus IV Epiphanes implemented oppressive policies against the Jewish people, including the prohibition of Jewish religious practices and the desecration of the Temple, which directly led to the Maccabean Revolt.
\item \textbf{Answer:} False
\\ \textit{Options:} A) True; B) False
\\ \textit{Note:} The statement is false because, while some interpretations of Hindu traditions may emphasize certain purity concepts, many Hindu texts and practices celebrate the divine feminine and recognize the vital and respected roles of women in society and spirituality.
\item \textbf{Answer:} False
\\ \textit{Options:} A) True; B) False
\\ \textit{Note:} Guru Nanak referred to the 'divine word' as 'Shabad' in his teachings, not 'Nam,' which is a different concept in Sikhism.
\end{enumerate}

\section*{Long Form Response}
\begin{enumerate}[label=\arabic*.]
\setcounter{enumi}{10}
\item \textbf{Answer:} The Digambara sect of Jainism holds specific beliefs regarding the role of women, particularly in the context of semi-renunciation. Digambaras believe that women cannot achieve full spiritual liberation without first being reborn as men, which reflects their view that women are at a disadvantage in attaining the highest spiritual goals. This contrasts with the Svetambara sect, which believes that women can attain liberation in their current form. Overall, the Digambara practices emphasize strict asceticism and often require male monks to renounce all possessions, while women are seen as less capable of following these rigorous paths. Thus, the differences in beliefs about women's spiritual potential highlight a significant divergence between the Digambara and other Jain groups.
\\ \textit{Note:} The answer correctly identifies the Digambara belief that women must be reborn as men to achieve full spiritual liberation, contrasting it with the Svetambara view that women can attain liberation in their current form, thus highlighting the differing beliefs and practices regarding semi-renunciation within Jainism.
\item \textbf{Answer:} The Daoist concept of wuwei, often translated as "non-action," emphasizes the importance of aligning with the natural flow of life rather than forcing actions or striving against the current. In Daoist philosophy, wuwei signifies a state of effortless action where individuals respond to situations in a harmonious and spontaneous manner, reflecting the balance of the Dao, or the fundamental principle of the universe. This approach encourages individuals to observe, adapt, and engage with the world around them without unnecessary struggle, leading to a more peaceful and fulfilling existence. By practicing wuwei, people can cultivate a deeper connection with nature and foster a sense of inner tranquility, ultimately promoting a more balanced way of living.
\\ \textit{Note:} The answer correctly identifies wuwei as a principle of non-action that promotes harmony with the natural flow of life, reflecting the core tenets of Daoist philosophy and its emphasis on spontaneous and adaptive interactions with the world.
\end{enumerate}

\vfill
\noindent\textit{End of Answer Key}
\end{document}